\documentclass{book}

\usepackage{listings}
\usepackage{xcolor}
\usepackage{colortbl}
\usepackage{hyperref}
\hypersetup{
    colorlinks=true,
    linkcolor=blue,
    citecolor=blue,
    urlcolor=blue
}
\usepackage{marginnote}


% vim:set errorformat="":
\usepackage{amsmath}
\usepackage{amsfonts}
\usepackage{amssymb}
\usepackage{xcolor}
\usepackage{colortbl}
\usepackage{hyperref}
\hypersetup{
    colorlinks=true,
    linkcolor=blue,
    citecolor=blue,
    urlcolor=blue
}
\usepackage{marginnote}
\usepackage{wrapfig}
\usepackage{graphicx}
\usepackage{lipsum}
\usepackage{array}
\usepackage{tabularx}
\usepackage{environ}
\usepackage{bussproofs}
\usepackage{exercise}

\renewcommand{\ExerciseHeader}{\medskip\noindent\textbf{\ExerciseName~\ExerciseHeaderNB.~\ExerciseHeaderDifficulty}\emph{\ExerciseHeaderTitle}\\~\\}
\renewcommand{\AnswerHeader}{\medskip\noindent{\textbf{Answer of \ExerciseName~\ExerciseHeaderNB}\\}}

\usepackage[skins,listings]{tcolorbox}
\usepackage{boolexpr,ifthen}
\definecolor{dkblue}{rgb}{0,0.1,0.5}
\definecolor{lightblue}{rgb}{0,0.5,0.5}
\definecolor{dkgreen}{rgb}{0,0.4,0}
\definecolor{dk2green}{rgb}{0.4,0,0}
\definecolor{dkviolet}{rgb}{0.6,0,0.8}
\definecolor{mantra}{rgb}{1.0,0.5,0}
\definecolor{gotcha}{rgb}{1.0,0.2,0}

\def\lstlanguagefiles{defManSSR.tex}
\lstset{language=SSR}
\newcommand{\mcbimpl}[1]{\lstinline/?$_{#1}$/}
\newcommand{\mcbimplm}[1]{\mbox{?}_{#1}}

\newcommand{\C}[1]{\mbox{\lstinlineþ#1þ}}
\let\L=\lstinline

% Highlights identifiers, |* id *| in underlined red, can we do better?
\lstset{moredelim=[is][\color{red}\bfseries\ttfamily\underbar]{|*}{*|}}

\newcommand\Chapter[2]{\chapter
  [#1\hfil\hbox{}\protect\linebreak{\itshape#2}]
  {#1\\ [2ex] {\Large \itshape #2}}
  \markboth{\MakeUppercase{\chaptername\ \thechapter.\ #1}}{}}

\newcolumntype{V}{>{\centering\arraybackslash} m{0\linewidth} }

\newcommand{\gotcha}[1]{
\begin{center}
\begin{tcolorbox}[colframe=gotcha,
	center title,
	title={\textcolor{white}{\textsf{\Large\textbf{Gotcha}}}}]
\begin{tabular}{lV}
\begin{minipage}{0.86\textwidth}
\vspace{1em}
\large
#1
\end{minipage}
&
\warnsign{1cm}
\end{tabular}
\end{tcolorbox}
\end{center}}

\newcommand{\mantra}[1]{
\begin{center}
\begin{tcolorbox}[colframe=mantra,
	center title,
	title={\textcolor{white}{\textsf{\Large\textbf{Mantra}}}}]
\begin{tabular}{lV}
\begin{minipage}{0.86\textwidth}
\vspace{1em}
\large
#1
\end{minipage}
&
\coqhead{1cm}
\end{tabular}
\end{tcolorbox}
\end{center}}

\newcommand{\coqhead}[1]{\includegraphics[width=#1]{../artwork/smallpantade}}
\newcommand{\warnsign}[1]{\includegraphics[width=#1]{../artwork/Ambox_warning_pn}}

\newcommand{\mcbpplevel}[1]{
\switch[#1=]
\case{1}
\case{2} \hspace{0.6cm} \coqhead{0.6cm}
\case{3} \coqhead{0.6cm} \coqhead{0.6cm}
\otherwise Please use \textbackslash mcbLEVEL
\endswitch
}

\newcommand{\mcbvarreset}{
\mcbLEARN{Please use \textbackslash mcbLEARN}
\mcbREQUIRE{Please use \textbackslash mcbREQUIRE}
\mcbPROVIDE{Please use \textbackslash mcbPROVIDE}
\mcbNOTES{}
\mcbLEVEL{4}}

\newcommand{\mcbLEARN}[1]{\def\mcbvarLEARN{#1}}
\newcommand{\mcbREQUIRE}[1]{\def\mcbvarREQUIRE{#1}}
\newcommand{\mcbPROVIDE}[1]{\def\mcbvarPROVIDE{#1}}
\newcounter{mcbvarLEVEL}
\newcommand{\mcbLEVEL}[1]{\setcounter{mcbvarLEVEL}{#1}}
\newcommand{\mcbNOTES}[1]{\def\mcbvarNOTES{#1}}
\newcounter{mcbvarTMP}

\newcommand{\mcbsection}[1]{
\setcounter{mcbvarTMP}{1}
\ifthenelse{\equal{}{\mcbvarNOTES}}{\setcounter{mcbvarTMP}{0}}{}
\section{#1}
\smash{\raisebox{0.7cm}{\hspace{0.85\textwidth}\mcbpplevel{\value{mcbvarLEVEL}}}}
\marginnote{\small\begin{tabular}{p{3.8cm}}\hline
\textbf{Learns} \mcbvarLEARN \\
\textbf{Requires} \mcbvarREQUIRE \\
\textbf{Provides} \mcbvarPROVIDE \\
\textbf{Level} \arabic{mcbvarLEVEL} \\
\switch[\value{mcbvarTMP}=]
\case{0}
\case{1} \textbf{Note} \mcbvarNOTES \\
\endswitch
\hline
\end{tabular}\\~\\
}\mcbvarreset{}
}

\newcommand{\mcbsubsection}[1]{
\setcounter{mcbvarTMP}{1}
\ifthenelse{\equal{}{\mcbvarNOTES}}{\setcounter{mcbvarTMP}{0}}{}
\subsection{#1}
\smash{\raisebox{0.7cm}{\hspace{0.85\textwidth}\mcbpplevel{\value{mcbvarLEVEL}}}}
\marginnote{\small\begin{tabular}{p{3.8cm}}\hline
\textbf{Learns} \mcbvarLEARN \\
\textbf{Requires} \mcbvarREQUIRE \\
\textbf{Provides} \mcbvarPROVIDE \\
\textbf{Level} \arabic{mcbvarLEVEL} \\
\switch[\value{mcbvarTMP}=]
\case{0}
\case{1} \textbf{Note} \mcbvarNOTES \\
\endswitch
\hline
\end{tabular}\\~\\
}\mcbvarreset{}
}

\mcbvarreset{}

\newbox{\devnull}
\NewEnviron{coqdef}[1]{}

\newcommand{\coqrun}[2]{}

\newtcblisting[auto counter]{coq}[2]{
	listing engine=listings,
	listing options={numbers=left,numberstyle=\tiny,
	                 aboveskip=0mm,belowskip=0mm,
		         xrightmargin=0mm,xleftmargin=0mm},
	%on line,
	before={\noindent},
	after={\hfill},
	listing only,
	fonttitle=\sffamily\bfseries,
	#2
}
\newtcblisting{coqout}[2]{
	listing engine=listings,
	listing options={aboveskip=0mm,belowskip=0mm,
	                 xrightmargin=0mm,xleftmargin=0mm},
	%on line,
	before={\noindent},
	after={\hfill},
	listing only,
	fonttitle=\sffamily\bfseries,
	#2
}


\title{The Book}
\author{Assia, Enrico, \ldots (you are welcome)}

\begin{document}
\maketitle
\tableofcontents{}

\setcounter{chapter}{-1}
\chapter{Introduction}

\emph{Mathematical Components}
is the name of a \Coq{} library covering a variety of
topics, from the theory of basic data structures (e.g., numbers, lists, finite
sets) to advanced results in various flavors of algebra. This library
constitutes the infrastructure for the formal proofs of the Four Color Theorem
and of the Odd Order Theorem.

The reason of existence of this book is to break down the barriers to entry.
While there are several books around covering the usage of the 
\Coq{}~\cite{BC04,SF,CPDT}
system and the theory it is based on~\cite[Chapter 4]{Coq:manual}\cite{paulinmohring:hal-01094195,hottbook}, 
the \mcbMC{} library is built
in an unconventional way.  As a consequence this book provides a non
standard presentation of \Coq{}, putting upfront the formalization choices
and the proof style that are the pillars of the library.

This books targets two classes of public.  On one hand newcomers, even the 
more mathematical inclined ones, find a soft introduction to the programming
language of \Coq{}, Gallina, and the Ssreflect proof language.
On the other hand accustomed Coq users find a substantial description
of the formalization style that made the \mcbMC{} library possible.

By no means this book pretends to be a complete description of \Coq{} or
Ssreflect: both tools already come with a comprehensive user 
manual~\cite{Coq:manual,ssrman}.
However, the reader is quickly put in the middle of a large library
of formalized concepts and given sufficient tools to prove non-trivial
mathematical results by reusing parts of the library.  
For example around page one-hundred the reader learns
how to show, formally, the infinitude of prime numbers, or the correctness of
the Euclidean's division algorithm, in a fewn lines of proof text.

%%%%%%%%%%%%%%%%%%%%%%%%%%%%%%%%%%%%%
\section{Structure of the book}

The book is divided into three parts.
The first two parts of the book are conceived to be read linearly,
while the third one can be read in any order.

\subsection{Part 1: Definitions and proofs}

This part introduces two languages and a formalization approach.

The first language is \emph{Gallina}, the mathematical foundation of the Coq
system.  For the expert reader: the type theory called \mcbCIC{}.  Such
language is both a programming and a specification language: it is used to
declare data, define programs, specify their behavior and represent their
correctness proofs.  Its most intriguing features are two: a rich system of
types, that can depend on other types or even values, and the notion of
computation, that is considered a first class citizen, i.e. a true form of
proof.

The second language, \mcbSSR{} (\emph{Ssreflect} for short), is used to write
proofs.  Its most characterizing features are two: it lets one write
concise and stable proof scripts enabling the development of large
libraries over a long time frame, and provides good support for the
formalization approach the \mcbMC{} library is based on, from which
its name is also taken.

The \emph{\mcbSSR{}} approach was initially conceived for the formal proof
of the Four Colors Theorem, and it later served as the pillar for
the \mcbMC{} library and the formal proof of the Odd Order Theorem.
Such approach takes major advantage from the symbolic computation
capabilities provided by Gallina.

The \Coq{} system provides a rich platform in which all this is
put together.  \Coq{} implements the Gallina language, and a checker for it.
It also provides a platform for the Ssreflect proof language and
many other facilities that are crucial for the \mcbMC{} library like
type inference.

\subsection{Part 2: Formalization techniques}

This part provides the tools to build a large library of formalized
mathematics.  In particular it presents a powerful form of automation
and a formalization technique that makes it possible to organize
concepts in a rational way and easily define new ones by linking them
to the already existing ones.

Automation is provided by \emph{programming type inference}.
The \Coq{} system provides a user-extensible type inference
algorithm.  It can be extended with declarative programs,
giving canonical solutions to otherwise unsolvable problems.
Such solutions typically involve notions and theorems that
are part of the \mcbMC{} library.  By programming type inference
one can hence teach \Coq{} the contents of the library.  The system
is then able to reconstruct non-trivial missing piece of information,
as an informed reader typically does when reading a mathematical text.

Formalized knowledge is organized by means of interfaces (in the spirit of
algebraic structures) and relations between them.  Type inference is programmed
to play the role of an librarian and recognize when an abstract theory has the
right to apply to a specific example.

Finally the rich language of \Coq{} lets one define new concepts
by refining existing ones, typically by gluing an object with
a proof of some extra property.  Type inference is programmed
to transport all the theory available on
the original concept to the derived one.

\subsection{Part 3: Mathematics in Mathematical Components}

This part provides a, unfinished by design, panoramic view of the
\mcbMC{} library.  It interleaves a catalog of formalized theories with
formalization choices that play a crucial role in the
definition of the formal object.

%%%%%%%%%%%%%%%%%%%%%%%%%%%%%%%%%%%%%%%%%%%%%%%%%%%%%
\section{Conventions}

Sections that are labelled with one \coqhead{0.3cm}, $(\star)$ in the table of
contents, are of medium complexity and are of interest to the reader willing
to extend the \mcbMC{} library.  Sections that are labelled with two
\coqhead{0.3cm}\coqhead{0.3cm}, $(\star\star)$ in the table of
contents, are of high complexity and are of interest to the reader willing
to understand the technical, implementation, details of the \mcbMC{} library.

Tricky details typically overlooked by beginners are signalled as follows:
\gotcha{Mind this detail.}

Advices one should keep in mind are signalled as follows:
\mantra{Remember this advice.}

\Coq{} code is in \C{typewriter} font and \C{(surrounded by parentheses)}
when it occurs the middle of the text.  Longer snippets are in boxes with line
numbers like the following one:

\begin{coq}{}{title=A sample snippet}
Example Gauss n : \sum_(0 <= i < n.+1) i = (n * n.+1) %/ 2.
Proof.
elim: n =>[|n IHn]; first by apply: big_nat1.
rewrite big_nat_recr //= IHn addnC -divnMDl //. 
by rewrite mulnS muln1 -addnA -mulSn -mulnS.
Qed.
\end{coq}

Code snippets are often accompanied by the goal status
printed by \Coq{}.

\begin{coqout}{}{title=Output from Coq after line 3}
n : nat
IHn : \sum_(0 <= i < n.+1) i = (n * n.+1) %/ 2
============================
\sum_(0 <= i < n.+2) i = (n.+1 * n.+2) %/ 2
\end{coqout}

% The Mathematical Components library was originally developed to
% support the formalization of the Odd Order theorem, a result that
% requires a wide selection of mathematical theories.
% 
% \section{challenges}
% 
% This introduction explains the motivation of the book, which is to
% present the methodologies we propose to build libraries of formalized
% mathematics \emph{that scale and can be reused}. Ideally, the content
% of the section is organized as a drawn-out comparison with the way
% mathematics on paper/blackboards are developed and written. Then the
% message is that in order to fulfill the technical needs, computer
% science techniques can be used and have worked well.
% 
% For instance, trying to formalize mathematical results from scratch in
% an empty context, without libraries and only with the tools proposed
% by the logic and the proof assistant is like trying to learn/invent
% mathematics without the help of the way mathematical concepts and
% notations have been shaped, structured, and denoted during the course
% of their (centuries long) study.
% In particular, this might lead the user to let the
% proof assistant impose a certain, usually more pedestrian proof, which
% we want to prevent as much as possible.
% (Another way to explain this could be to imagine how to explain some math
% to a kid that received no education, or trying to imagine that kid
% inventing new math... the cultural backround is important to be up to
% speed. Maybe the aim of mathcomp is to implement that background at the
% quality of today's math).
% 
% There are several important issues that have to be solved in order to
% come up with a library that has a chance to be reusable. The most
% pervasive one is the line drawn between what is trivial and what is
% logged in the proofs (for humans). Some are obvious (reflexivity of
% relations, ...) but in general this should be analyzed by
% \marginnote{Lost in the ...}
% ... learning the maths and reading the literature in the domain of
% interest. This appreciation might also be depending on the knowledge
% the reader has: in 1st year texts you make dozen of proofs that
% certain sets are equipped with a structure of group, vector space or
% whatever, but very soon you're supposed to \emph{see} that it is
% obviously the case. Emphasizing this point and providing a methodology to
% implement this in a systematic way might be one of the original points
% of this document. 
% 
% One of the great achievements attributed to Bourbaki is the
% generalization of the so-called axiomatic method, that promotes the
% design and use of abstract mathematical structures that factor theory
% and notations. This not only makes maths more readable, but also more
% understandable and you need to play with the concepts for some times
% to come up with the right abstractions. Incidentally, notations are
% not only overloading of symbols, but also carry inference of
% properties. We should be able to do this and find a way to capture this
% inference, which is in fact a (Prolog like) program run by the trained
% reader. And in general we would like the user to provide as much, but
% no more information in his statement/proof as he would in a proof,
% avoiding redundant information (where redundant covers inferable).
% 
% Another kind of implicit information is behind reasoning patterns that
% should also be modeled in a convenient way for the user. For instance,
% \emph{mutatis mutandis}, wlog, etc. These rely on the fact that the
% reader is able both to infer the mathematical statement involved in a
% cut formula \marginnote{LR: is it reasonable to talk about cut formula here?}
%  and to run rather easily the simple piece of proof that
% justifies it. This is modeled by techniques based on formula generation
% by subterm selection. This also overcomes an unwanted behavior that
% accessing to a subterm depends on the depth at which it occurs in the
% formula (usual bureaucracy in naive formalization).
% 
% The techniques that make this possible are adapted from standard
% methods in computer science and programming language research. We stand on the
% shoulders of a foundational system based on a type theory, which
% promotes functions as the initial concept of the formalism (as opposed
% to sets). In particular computations plays a privileged role there,
% allowing to model a notion of ``similar'' up to (implicit)
% computations. Things that are equal up to computations can be
% substituted one by the other in a transparent way. This is central to
% the way things are modeled and at the heart of (small) scale
% reflection. Computer science engineering techniques are also useful
% to organize the infrastructure content that make it possible to work with
% some comfort in upper layers of the library, in order to configure
% the content so that the machine can use it (order of arguments,
% implicit status, naming policy...).
% 
% 
% \section{challenges faced and tools adopted}
% (tools in a broad sense, the logic is a tool, coq is a tool, the plugin
% is a tool, the ssr style is a tool,...)
% 
% Challenges:
% \begin{itemize}
% \item large body (scale up), make proofs small and robust.
% 	We need to say that we do use "deterministic automation".
% 	Use Laurent's data on de bruijn factor.
% \item model the use of math notations, their role in proofs, model proofs (also
% 	it is about reasoning, not just computations). Another way to say that
% 	is: model Bourbaki (rationalization of Math via structure/interfaces)
% 	but not the first book (set theory) that is replaced by CIC (link with
% 	section computational thinking).
% \end{itemize}
% 
% We build on Coq and an extension.  The main tools follow (in random order):
% 
% \section{Computational Thinking}\label{ch:compthink}
% 
% This section should motivate the activity of formalizing mathematics
% with the {\bf Coq} proof assistant, emphasizing its computing skills,
% and as opposed to other foundations like HOL. However the challenge is
% to keep mathematicians as the privileged target, while motivating the
% ssr approach to a CS oriented reader.
% 
% See \verb+../coq/ch0.v+.
% 
% Aim of the chapter:
% \begin{itemize}
% \item should sound natural and easy to a CS person (but with the ssr twist)
% \item should sound different but well motivated to a Coq user (do show, maybe in
%   the exercises, that leqn is 100 times better than "Inductive le").  Try to
%   reproduce the shock we had the first time we used Boolean predicates.  It may
%   help to compare, in the *advanced* section, the approach with the standard
%   one, so that one sees two proof scripts in the same page.
% \end{itemize}
% 
% \section{logic programming ... type inference}
% to model proof search, give a meaning to notations, teach coq the
% work an informed reader does (contextualizing otherwise ambiguous
% notations, knowledge of interface/instance of algebraic structures).
% 
% relation with proof search: no "blind" proof search (easy, ad-hoc, pervasive
% v.s. advanced, generalistic, potentially expensive and unstable).
% 
% \section{automation in tactics}
% the main points:
% \begin{itemize}
% \item 1/3 is rewrite, term selection/search (one does not need to reach a sub
% 	formula as a goal in order to make progress for example, no monkey
% 	puzzle as in GG terminology).
% \item Create a formula without writing it: some advanced forms of forward
% 	reasoning tactics to deal with symmetries, generalizations (boils down
% 	to syntesize the cut formula out of the minimum possible user input, as
% 	in a text where one says "similarly to that, we can also prove that".
% 	(this is not very pervasive, dunno if it is worth putting it here in
% 	this chapter).  Also elim does that. (technically also rewrite, but we
% 	may want to separate things)
% \end{itemize}
% This section is were one talks about the plugin, and some of the main design
% points of tactics: compositional (a language, not a list of commands),
% predictable (documented!!!), finally compact (symbols for uninteresting steps).
% 
% \section{discipline}
% MAKE an howto out of that.
% Maybe one should also add a few notes on the style in scripts? like:
% \begin{itemize}
% \item you must be able to model (at least) 1 proof step in a sentence (line),
%   e.g. "rewrite preparegoal dostep ?cleanup."
% \item uninteresting/recurrent lemmas/steps should be small (short names, easy to
%   gray out)
% \item lemma statements are designed, not just written, having in mind their use
%       (forward, backward, implicit argumets, arguments order)
%       and the class of trivial hypotheses since an extra hyp that is proable
%       triviality (via //, hint resolve, cnaonical) is for free. E.g.
%       "x \\is a toto", "0 <= n", ...
% \item also not every possible lemma, but a few that combine well
% \item proofs/definitions are reworked many times, why (understand recurrent
%   proof schemas, compact, factor, make more stable/robust) and what is needed
%   (like meaningful names, clear structure)
% \end{itemize}
% 
% \section{trivial=implicit (for a trained mathematician)}
% The idea is to try to identify what is trivial (mathematically speaking)
% and be sure you can model it as such:
% \begin{itemize}
% \item (level basic) make explicit the trivialities of each theory (what one
% 	expects to be proved by //). 
% \item (level advanced) when you do new stuff, you must decide what is
% 	trivial/implicitly proved.
% \item (hard) which technique to make Coq prove it automatically (hint resolve,
% 	canon, comput... in the type)
% \end{itemize}
% 
% % This may also be another way present the whole chapter
% 
% \mantra{(basic) if you see toto=false you should perform the case analysys via
% fooP
% }




\part{The art of formalizing}

\chapter{Logics}

From calculability to proofs, hence the CC, and the fact that
reasoning principles without a computational content become axioms.

This is a non technical chapter and message should be:
\begin{itemize}
\item instantiation of a universal statement is application (also the pair)
\item Excluded middle is not available by default (choice?)
\item Conversion as a pervasive indistinguishably, what inside
  (beta, definition unfolding,...)
\item Dependent types: eq, sigma (which example?)
\end{itemize}

One options is: avoid relating type theory and other logics. We say:
we have a formal game where the basic elements are programs/functions
that come with types to avoid confusion. full stop. (no relation with
proof theory, set theory). maybe mention that roots are in calculability (hence
the choice to pick functions as primitive and not sets). This is lucky because
(computable) functions are today executable by a computer.  Still not all
concepts are "computable" hence some principles are problematic: EM,.... we
mainly stay in the lucky fragment (again no propaganda on intuitionistic logic,
constructive math; just a mention).

\chapter{Programming}

Presentation of inductive data structures, recursive programs on these
data.
Bool, Boolean connectives, Boolean reflection (cf
ch. \ref{ch:compthink}), views.
Examples of equality tests ($==$, with a forward reference to
explain the magic if needed), operations on sequences, nats,
exercises on prime, div.

\chapter{Proofs}

Where one learns to do proofs.
Boolean reflection in practice, views, discussion on the definition of leq,
proofs on things defined in the
previous chapter, associated tactics, exercises on prime, div,
binomial, etc.

spec? A new vernacular to declare specs without typing coinductive and
by writing explicitly the equations.

\chapter{Type inference/Genericity}

put here a minimalistic presentation of records.

Implicit arguments (Is it possible to survive without implicit
argument up to this point? We should probably use them plus forward
references to here), canonical structures. Combining with
notations. Now the real $==$. 

\chapter{Mixing data and proofs}
Boolean sigma types, records, coercions, UIP. Examples: ordinals,
tuples (not their use which requires CS). This chapter might come
after some presentation of basic canonical structures, i.e. reorganize
the content between this chapter and the next.

example of tuples here? better ordinals?

GG: makes many points here (not fully understood by Enrico):
\begin{itemize}
\item ordinals/tuples are easy to use but hard to build
\item it is a tradeoff, but is not clear if we can give hints on when
	a specific datatype like ordinals is better that unpackaged
	stuff.
\end{itemize}

UIP is advanced, the basic user should just be told that putting
bools inside a record is just fine.

\begin{itemize}
	\item record (flat)
	
\end{itemize}


\chapter{Hierarchy}
Packaging records, the bigop hierarchy.
Scaling with packed classes and mixins, to the ssralg
hierarchy. Presentation of the content of ssralg in terms of structures
and of the theory? Should the latter be a separate chapter.

Maybe a plugin for a new vernacular to script the creation/declaration
of structures/instances so that the level basic can touch the argument
easily.

Explain what the abstraction barrier is (like unfolding a GRing projection)

\gotcha{if you see GRing.toto then you broke an abstraction barrier}

\mantra{
	if you have a proof in mind, don't let the system drive you
        to another, less clean and abstract, proof.
}

Declaring an instance is hard... we need to document the multuiple processes 
for each structure in each hierarchy and possibily make a program out of it.

\chapter{Larger scale reflection (out of place)}
Four colours, decomposition in primes, example in peterfalvi. In
particular example where a case analysis does not follow the
constructors of an existing inductive type: then craft an ad hoc one
to hint the proof. This topic is probably at a wrong place.

\part{Mathematics in Mathematical Components}

\section{STYLE of these chapters}
These chapters hopefully in the following style:
\begin{itemize}
\item first math mode to review the (not so standard) definitions and
\item then touch some real proof script to show why/how the definition
  make it possible to model the proofs (as they are in math)
\item use the (few) examples to illustrate a cool proof strategy if
  any (not necessarily typical of the math subject, but that happens to
  be there, like in OOTHM paper: circular leq, symmetries, ad hoc
  decision procedure, ...)
\item try hard to show how CIC helps (which feature: HO, computation, dependent
	types \ldots), so that at the end we can sum up and make a synthesis of
	all that.
\end{itemize}

Maybe it is simpler to do it in 2 steps:
1) in this second part one identifies where CIC or the SSR style plays a
   crucial role
2) then we advertise these use cases in the first part to motivate
   the techniques, the complexity of the logic...

\chapter{Numbers}

What are the numbers available in the library? How to use them, casts
between types... Non trivial point in the formalisation:
axiomatization, order is defined from norm, partial orders (in
particular for complex numbers).


\chapter{Polynomials, Linear algebra}

2-stage presentation: interface plus explicit. Expansion of
Georges'ITP paper. Here is one of the main application of the choice
operator (complement a base).

\chapter{Quotients}

\chapter{Finite group theory}

Data-structures, how to craft the set of variants of a same theorem to
make the formalization handy. Permutations. Presentations?

\chapter{Representation theory, Character theory}

As an example of application, in particular of the linear algebra
theory.

\chapter{Galois Theory}


\chapter{Real closed fields}

\chapter{Algebraic closure}

\part{Annexes}

\chapter{What is done where?}

\chapter{How tos}
suggestions:
\begin{itemize}
\item Get more information when you do not understand the error
  message
\item Search in the library
\item Canonical structures: define a new instance
\item Canonical structures: add a new structure
\item Give a relevant name to the lemma I just stated
\item Forbid unwanted expansions
\item Choose a notation (what not to do...)
\item Compute ``for real'' (Natbin, Coqeal)
\item MathComp script style
\end{itemize}
\chapter{Naming conventions}

\chapter{Index of notations}

\end{document}
